%%%%%%%%%%%%%%%%%%%%%%%%%%%%%%%%%%%%%%%%%%%%%%%%%%%%%%%%%%%%%%%%%
                % Deterministic transformer %
%%%%%%%%%%%%%%%%%%%%%%%%%%%%%%%%%%%%%%%%%%%%%%%%%%%%%%%%%%%%%%%%%

We call a transformer deterministic when the token selector is the one taking $\argmax$ over the distribution.
We say that a family of deterministic transformers $\{\TF_n\}_{n \in \mathbb{N}}$ recognize a language $\gL$ when: 
\[
\alpha \in \gL \iff \TF_{|\alpha|}^*(\alpha) = \qaccept
\]

Given two functions $T(n)$ and $d(n)$ we define the complexity class $\CoT[T(n), d(n)]$.
Informally, a language $\gL$ is in $\CoT[T(n), d(n)]$ iff there is a deterministic family of transformers with constant depth, embedding size $d(n)$ and budget $T(n)$ that recognizes it.
Formally, we say a language  $\gL$ is in $\CoT[T(n), d(n)]$ iff there is an integer $L$ and two functions $T'(n) = O(T(n)), d'(n) = O(d(n))$, such that for every positive integer $n$, there is an $L$-layer deterministic transformer, denoted by $\TF_n$ with embedding size $d'(n)$, that can output $\gL(\alpha)$ given any input $\alpha$ in $\Sigma^n$, using $T'(n)$ steps of chain of thought.


